\documentclass[letterpaper]{article} 
\usepackage[utf8]{inputenc}
\linespread{0.85}
\usepackage[T1]{fontenc}
\usepackage{amsmath}
\usepackage{amsfonts}
\usepackage{amssymb}
\usepackage{array}
\usepackage{fontspec}
\usepackage{booktabs}
\usepackage{hyperref}
\usepackage[version=4]{mhchem}
\usepackage{stmaryrd}
\usepackage[dvipsnames]{xcolor}
\colorlet{LightRubineRed}{RubineRed!70}
\colorlet{Mycolor1}{green!10!orange}
\definecolor{Mycolor2}{HTML}{00F9DE}
\usepackage{graphicx}
\usepackage{amsmath}
\usepackage{graphicx}
\usepackage{capt-of}
\usepackage{lipsum}
\usepackage{fancyvrb}
\usepackage{tabularx}
\usepackage{listings}
\usepackage[export]{adjustbox}
\graphicspath{ {./images/} }
\usepackage[utf8]{inputenc}
\usepackage[english]{babel}
\usepackage{float}
\usepackage{lipsum}
\usepackage{graphicx}
\usepackage{float}
\usepackage[margin=0.7in]{geometry}
\usepackage{amsmath}
\usepackage{graphicx}
\usepackage{capt-of}
\usepackage{tcolorbox}
\usepackage{lipsum}
\usepackage{graphicx}
\usepackage{float}
\usepackage{listings}
\definecolor{deepred}{RGB}{139,0,0}
\usepackage{hyperref} 
\usepackage{xcolor} % For custom colors
\lstset{
	language=Python,                % Choose the language (e.g., Python, C, R)
	basicstyle=\ttfamily\small, % Font size and type
	keywordstyle=\color{blue},  % Keywords color
	commentstyle=\color{gray},  % Comments color
	stringstyle=\color{red},    % String color
	numbers=left,               % Line numbers
	numberstyle=\tiny\color{gray}, % Line number style
	stepnumber=1,               % Numbering step
	breaklines=true,            % Auto line break
	backgroundcolor=\color{black!5}, % Light gray background
	frame=single,               % Frame around the code
}
\usepackage{float}
\usepackage[]{amsthm} %lets us use \begin{proof}
	\usepackage[]{amssymb} %gives us the character \varnothing
	
	\title{Assignment 4, MATH 4155}
	\author{Zongyi Liu}
	\date{Sun, Mar 22, 2026}
	\begin{document}
		\maketitle
		
		\section{Question 1}
		
		{\color{deepred}\fontspec{Georgia} HW8 Q1 Optimal Stopping Time}
		
		Sometimes, bad news is better than no news at all. Let us consider the simple, symmetric random walk $\left\{S_{n}\right\}_{n \in \mathbb{N}_{0}}$ started at $S_{0}=0$, as a model for the fluctuations of the price of some financial asset in our portfolio (which may well take negative values, if we happen to owe money).
		
		Assume that Martha has been told that at some time $N$ with $N \in\{7,9,11, \cdots\}$ the asset is going to be worth $S_{N}=-1$. Show that by selling the asset at time
		
		$$
		T=\min \left\{n \in \mathbb{N}_{0} \mid S_{n}=1\right\} \wedge N
		$$
		
		(that is, by selling the first moment the asset is worth $\$ 1$ if there is such a moment before time $N$, and selling at time $N$ otherwise), Martha can secure a positive expected resell value
		
		$$
		\mathbb{E}\left(S_{T} \mid S_{N}=-1\right)>0 ;
		$$
		
		and this, despite the fact that initially her asset is worthless ( $S_{0}=0$ ), and despite the fact that she even knows things are going to get 'bad' in the end $\left(S_{N}=-1\right) .^{1}$\\
		(\emph{Hint}: Recall $\mathbb{E}(X \mid A)=\mathbb{E}\left(X \mathbf{1}_{A}\right) / \mathbb{P}(A)$ and use the reflection principle to rewrite the expression you get when you apply this formula to $S_{T}(\omega)=S_{T(\omega)}(\omega)$. This will allow you to reduce the claim $\mathbb{E}\left(S_{T} \mid S_{N}=-1\right)>0$ to an inequality for the distribution of the random variable $S_{N}$, which you can then verify easily.)
		
		
		\textbf{Answer}
		
		\clearpage
		
		\section{Question 2}
		
		{\color{deepred}\fontspec{Georgia} HW8 Q2 Gumbel Domain of Attraction}
		
		\begin{itemize}
			\item[(i)] Let $X_1,X_2,X_3,\cdots$ be independent random variables with common exponential distribution with parameter $\lambda>0$. What can you say about the distribution of the random variable
			\[
			Y_n=\lambda\left(\max_{1\le j\le n} X_j\right)-\log n,
			\qquad \text{as } n\to\infty?
			\]
			
			\item[(ii)] Let $X_1,X_2,X_3,\cdots$ be independent random variables with common ``logistic'' distribution
			\[
			\mathbb{P}(X_j\le x)=F(x)=\frac{1}{1+e^{-x}},\qquad x\in\mathbb{R}.
			\]
			What can you say about the distribution of the random variable
			\[
			Y_n=\max_{1\le j\le n} X_j-\log n,
			\qquad \text{as } n\to\infty?
			\]
		\end{itemize}
		
		\textbf{Answer}
		
			
		
		\clearpage
		
		\section{Question 3}
		
		{\color{deepred}\fontspec{Georgia} HW8 Q3 Riemann-Zeta Calculus} 
		
		Show that
		
		$$
		\lim _{n \rightarrow \infty} e^{-n} \sum_{k=0}^{n} \frac{n^{k}}{k!}=\frac{1}{2} .
		$$
		
		\textbf{Answer} 
		
		
		
		\clearpage
		
		\section{Question 4}
		
		
		{\color{deepred}\fontspec{Georgia} HW8 Q7 Walsh 5.11}
		
		Let $X_1,X_2,\ldots$ be independent random variables with mean zero and finite variance. Put $S_n=X_1+\cdots+X_n$.
		\begin{enumerate}
			\item[(a)] Let $\varepsilon>0$ and set
			\[
			\Lambda=\{\omega:\max_{1\le k\le n}|S_k(\omega)|\ge \varepsilon\}.
			\]
			Define
			\[
			\Lambda_k=\{\omega:\max_{1\le j\le k-1}|S_j(\omega)|<\varepsilon,\ |S_k(\omega)|\ge \varepsilon\},
			\qquad k=1,\ldots,n.
			\]
			Show that the $\Lambda_k$ are disjoint and $\Lambda=\bigcup_{k=1}^n \Lambda_k$.
			
			\item[(b)] Show that $S_n-S_k$ is independent of $S_k$.
			
			\item[(c)] Show that
			\[
			\int_{\Lambda_k} S_n^2\,dP \ge \int_{\Lambda_k} S_k^2\,dP.
			\]
			[Hint: $S_n=S_k+(S_n-S_k)$.]
			
			\item[(d)] Deduce that
			\[
			E\{S_n^2\}\ge \sum_{k=1}^n \int_{\Lambda_k} X_k^2\,dP
			\ge \varepsilon^2 P\{\Lambda\}.
			\]
			
			\item[(e)] Conclude {Kolmogorov's Inequality}:
			\begin{equation}
				P\left\{\max_{1\le k\le n}|S_k|\ge \varepsilon\right\}
				\le \frac{1}{\varepsilon^2} E\{S_n^2\}.
				\tag{5.4}
			\end{equation}
		\end{enumerate}
		
		Note that if $\max_{1\le k\le n}|S_k|$ is replaced by $|S_n|$ itself, this is just Chebyshev's inequality. So Kolmogorov's inequality is a significant extension of Chebyshev's inequality. Inequalities like this will reappear in Chapter~9.
		
		
		\textbf{Answer} 
		
		
		
		\clearpage
		
		\section{Question 5}
		
		{\color{deepred}\fontspec{Georgia} HW8 Q8}
		
		\subsection{Part 1}
		{\color{deepred}\fontspec{Georgia}  Walsh 5.12 Convergence of Series}
		
		Let $X_1,X_2,\ldots$ be independent random variables with mean zero and finite variance. Put $S_n=X_1+\cdots+X_n$. Show that if $\sum_{n=1}^{\infty} E\{X_n^2\}<\infty$, then $\sum_{n=1}^{\infty} X_n$ converges a.s.
		
		
		
		(\emph{Hint}: Use Kolmogorov's inequality (Exercise 5.11).)
		
		
		\textbf{Answer}
		
		\clearpage
		
		\subsection{Part 2}
		{\color{deepred}\fontspec{Georgia}  Walsh 5.13 Random Signs Problem}
		
		Let $(a_n)$ be a sequence of real numbers. Choose the sign, positive or negative, of $a_n$ at random. Does $\sum_n \pm a_n$ converge? To answer this, let $X_1,X_2,\ldots$ be i.i.d.\ with $P\{X_n=1\}=P\{X_n=-1\}=1/2$. Let $(a_n)$ be a sequence of real numbers. Show that if $\sum_n a_n^2<\infty$, then $\sum_n a_n X_n$ converges a.s. Note: The random series actually converges if and only if the sum of squares converges, but the converse is slightly harder. The definitive result on the convergence of random series is Kolmogorov's Three Series Theorem.
		
		
		\textbf{Answer} 
		
		
		
		\clearpage
		
		
		
		
	\end{document}
